\documentclass[12pt]{article}
\usepackage[utf8]{inputenc}
\usepackage{amsmath,amssymb,amsthm}
\usepackage{hyperref}
\usepackage{graphicx}
\usepackage{booktabs}
\usepackage{listings}
\usepackage[margin=1in]{geometry}

\title{A Twelve-Stage Verification Pipeline for\\Computational Physics Research}

\author{John Roehm}
\date{Draft --- April 2026}

\begin{document}

\maketitle

\begin{abstract}
We describe a twelve-stage pipeline for conducting computational physics
research with continuous machine verification. The pipeline integrates
Python numerics, Lean~4 formal proofs, and the Aristotle automated theorem
prover into a single workflow where every formula has a formal theorem,
every computed quantity has verified bounds, and every paper claim traces
to computation. Applied to a research program spanning dissipative effective
field theory, analog Hawking radiation, emergent gravity, quantum groups,
and topological quantum field theory, the pipeline has produced
% Auto-generated by update_counts.py — 2026-04-24T19:44:03.140400
% Do not edit manually. Run: uv run python scripts/update_counts.py
\newcommand{\totaltheorems}{3993}
\newcommand{\substantivetheorems}{3882}
\newcommand{\placeholdertheorems}{111}
\newcommand{\axiomcount}{1}
\newcommand{\sorrycount}{0}
\newcommand{\sorrypercent}{0.0}
\newcommand{\leanmodules}{167}
\newcommand{\leandefinitions}{3297}
\newcommand{\pythonmodules}{68}
\newcommand{\testfiles}{59}
\newcommand{\totaltests}{2836}
\newcommand{\figurecount}{111}
\newcommand{\notebookcount}{54}
\newcommand{\papercount}{19}
\newcommand{\aristotleproved}{322}
\newcommand{\aristotleruns}{44}
% --- Paper 18 / Phase 5t per-section counts (FermiHubbardDimer.lean) ---
\newcommand{\fhdTotal}{143}
\newcommand{\fhdWaveTwo}{13}
\newcommand{\fhdWaveThree}{6}
\newcommand{\fhdWaveFour}{22}
\newcommand{\fhdWaveFive}{22}
\newcommand{\fhdWaveSixABC}{38}
\newcommand{\fhdWaveSixExtra}{15}
\newcommand{\fhdWaveSeven}{19}
\newcommand{\fhdWaveEight}{8}
\totaltheorems{} machine-checked theorems
across % Auto-generated by update_counts.py — 2026-04-24T19:44:03.140400
% Do not edit manually. Run: uv run python scripts/update_counts.py
\newcommand{\totaltheorems}{3993}
\newcommand{\substantivetheorems}{3882}
\newcommand{\placeholdertheorems}{111}
\newcommand{\axiomcount}{1}
\newcommand{\sorrycount}{0}
\newcommand{\sorrypercent}{0.0}
\newcommand{\leanmodules}{167}
\newcommand{\leandefinitions}{3297}
\newcommand{\pythonmodules}{68}
\newcommand{\testfiles}{59}
\newcommand{\totaltests}{2836}
\newcommand{\figurecount}{111}
\newcommand{\notebookcount}{54}
\newcommand{\papercount}{19}
\newcommand{\aristotleproved}{322}
\newcommand{\aristotleruns}{44}
% --- Paper 18 / Phase 5t per-section counts (FermiHubbardDimer.lean) ---
\newcommand{\fhdTotal}{143}
\newcommand{\fhdWaveTwo}{13}
\newcommand{\fhdWaveThree}{6}
\newcommand{\fhdWaveFour}{22}
\newcommand{\fhdWaveFive}{22}
\newcommand{\fhdWaveSixABC}{38}
\newcommand{\fhdWaveSixExtra}{15}
\newcommand{\fhdWaveSeven}{19}
\newcommand{\fhdWaveEight}{8}
\totalmodules{} Lean modules with
% Auto-generated by update_counts.py — 2026-04-24T19:44:03.140400
% Do not edit manually. Run: uv run python scripts/update_counts.py
\newcommand{\totaltheorems}{3993}
\newcommand{\substantivetheorems}{3882}
\newcommand{\placeholdertheorems}{111}
\newcommand{\axiomcount}{1}
\newcommand{\sorrycount}{0}
\newcommand{\sorrypercent}{0.0}
\newcommand{\leanmodules}{167}
\newcommand{\leandefinitions}{3297}
\newcommand{\pythonmodules}{68}
\newcommand{\testfiles}{59}
\newcommand{\totaltests}{2836}
\newcommand{\figurecount}{111}
\newcommand{\notebookcount}{54}
\newcommand{\papercount}{19}
\newcommand{\aristotleproved}{322}
\newcommand{\aristotleruns}{44}
% --- Paper 18 / Phase 5t per-section counts (FermiHubbardDimer.lean) ---
\newcommand{\fhdTotal}{143}
\newcommand{\fhdWaveTwo}{13}
\newcommand{\fhdWaveThree}{6}
\newcommand{\fhdWaveFour}{22}
\newcommand{\fhdWaveFive}{22}
\newcommand{\fhdWaveSixABC}{38}
\newcommand{\fhdWaveSixExtra}{15}
\newcommand{\fhdWaveSeven}{19}
\newcommand{\fhdWaveEight}{8}
\sorrycount{} remaining gaps, while
maintaining 16 automated validation checks and full parameter provenance.
We present the pipeline architecture, its ten invariants, the automated
theorem prover integration that compresses months of manual proof work
to hours, and the claims-reviewer system that catches inconsistencies
between papers and computation. The methodology is domain-agnostic and
applicable to any research program combining numerical computation with
formal verification.
\end{abstract}

% ============================================================
\section{Introduction}
\label{sec:intro}

% Motivation: computational physics has a reproducibility crisis
% - Numerical results without verified bounds
% - Paper claims that diverge from computation over time
% - No machine-checkable connection between theory and code
%
% Our solution: a pipeline that makes verification continuous, not post-hoc
% Applied to SK-EFT Hawking: BEC analog Hawking radiation + quantum groups + MTCs

% ============================================================
\section{Pipeline Architecture}
\label{sec:pipeline}

% The 12 stages:
% 1. Constants & Parameters (single source of truth)
% 2. Formulas (canonical, with Lean refs)
% 3. Lean Theorems (sorry stubs)
% 4. Aristotle (automated proving)
% 5. Lean Build Verification (zero sorry gate)
% 6. Python Tests (correctness + physical bounds)
% 7. Cross-Layer Validation (16 checks)
% 8. Visualizations (publication-quality, from validated data)
% 9. Figure Review (LLM visual inspection)
% 10. Paper Draft (every claim traced)
% 11. Notebooks (reproducible narratives)
% 12. Document Sync (automated counts)

% ============================================================
\section{Pipeline Invariants}
\label{sec:invariants}

% The 10 invariants that must hold at ALL times:
% 1. formulas.py is canonical
% 2. constants.py is canonical
% 3. visualizations.py is canonical
% 4. Every formula has a Lean theorem (zero sorry)
% 5. Every computed quantity has bounds
% 6. Every paper claim traces to computation
% 7. Narrative derives from data
% 8. Every parameter has verified provenance
% 9. Placeholder theorems are non-load-bearing
% 10. No heartbeat overrides

% ============================================================
\section{Three-Layer Verification}
\label{sec:three_layer}

% Layer 1: Python numerics (fast iteration, physical intuition)
% Layer 2: Lean 4 formal proofs (mathematical certainty)
% Layer 3: Aristotle automated theorem prover (fills sorry gaps)
%
% The interaction: Python discovers → Lean formalizes → Aristotle proves
% Example: tetrad gap equation G_c = 8pi^2/(N_f Lambda^2)

% ============================================================
\section{Automated Theorem Prover Integration}
\label{sec:aristotle}

% Aristotle: commercial Lean 4 proving service
% Integration pattern:
% - Priority batching (P1/P2/P3)
% - PROVIDED SOLUTION hints in docstrings
% - One job at a time (sends entire project)
% - Manual review before integration
% - Quality verification (exercises hypotheses?)
%
% Results: 307 theorems proved automatically
% Compression: "multi-month estimates" → hours/days
% Failure mode: OUT_OF_BUDGET (need to split or hint better)

% ============================================================
\section{Validation Infrastructure}
\label{sec:validation}

% 16 automated checks in validate.py:
% 1. Formula-Lean cross-reference
% 2. Numerical parameter verification
% 3. Mathematical identities
% 4-6. Paper/notebook/lean consistency
% 7-11. Cross-path, bounds, provenance
% 12-16. Graph integrity, parameter provenance
%
% Claims-reviewer: LLM agent that reads papers against computation
% Figure-reviewer: LLM agent that reviews visualizations
% Provenance dashboard: human verification interface

% ============================================================
\section{Automated Counts System}
\label{sec:counts}

% Problem: manual count synchronization across 10+ documents
% Solution: ExtractDeps.lean → lean_deps.json → update_counts.py → counts.json
% Single source of truth for all project metrics
% LaTeX macros (\totaltheorems, \sorrycount, etc.) auto-generated

% ============================================================
\section{Case Study: From Discovery to Verified Theorem}
\label{sec:case_study}

% Walk through a specific example end-to-end:
% The pentagon equation convention bug
% 1. Aristotle "proved" a wrong pentagon (vacuous)
% 2. Brute-force search identified convention mismatch
% 3. Deep research provided correct Kitaev convention
% 4. native_decide verified 243 cases in seconds
% 5. Claims-reviewer caught stale paper references
%
% Lesson: the pipeline CAUGHT a real bug that would have propagated

% ============================================================
\section{Decidable Number Fields}
\label{sec:number_fields}

% The technique for verified algebraic computation
% Q[x]/(f(x)) as rational tuples with DecidableEq
% native_decide for pentagon (243 cases), hexagon (6 eqs), S-matrix unitarity
% Trust trade-off: compiler (~30K lines) vs kernel-only
% Performance: seconds vs hours for equivalent norm_num proofs

% ============================================================
\section{Scaling and Applicability}
\label{sec:scaling}

% Current scale: 2232 theorems, 94 modules, 14 papers, 42 notebooks
% Growth rate: ~200 theorems per session with the pipeline
% Bottlenecks: Aristotle budget, heartbeat limits, number field consolidation
%
% Applicability to other domains:
% - Pharmaceutical research (FDA SaMD algorithmovigilance)
% - Financial modeling (verified risk calculations)
% - Climate science (verified model predictions)
% - Any domain combining computation with formal claims

% ============================================================
\section{Related Work}
\label{sec:related}

% PhysLean / HepLean: quantum harmonic oscillator, Higgs potential
% Verified quantum computing: SQIR/VOQC, Lean-QuantumInfo
% Reproducibility tools: Jupyter, papermill, DVC
% Formal methods in industry: CompCert, seL4, CertiKOS
% How this differs: continuous verification during research, not post-hoc

% ============================================================
\section{Conclusion}
\label{sec:conclusion}

% The pipeline transforms research methodology:
% - Bugs caught in real time (pentagon convention)
% - Claims verified against computation automatically
% - Aristotle compresses proof work by orders of magnitude
% - The 10 invariants prevent drift between papers and code
%
% Open-source, domain-agnostic, immediately applicable

\bibliographystyle{plain}
\bibliography{refs}

\end{document}
